\chapter{Behavioral State Model}

The Behavioral Engine for Demand (BED) models individual behavior using a
low-dimensional, interpretable state representation that captures how users
respond to incentives over time. This chapter formalizes the behavioral state
variables, transition dynamics, and event probability model that form the core
of BED v3. The behavioral model integrates insights from habit formation,
recency effects, incentive decay, and heterogeneous responsiveness.

\section{Overview}

BED represents each user with a latent behavioral state:


\[
s_t = (r_t, h_t, \rho_t, \Delta_t, I_{t-1}),
\]


where:
\begin{itemize}
    \item $r_t$ = recency (time since last event),
    \item $h_t$ = habit strength,
    \item $\rho_t$ = responsiveness to incentives,
    \item $\Delta_t$ = time since last incentive,
    \item $I_{t-1}$ = magnitude of the last incentive delivered.
\end{itemize}

This state is a sufficient statistic for predicting event probability and
determining optimal incentive timing, consistent with the BED\_MIT.pdf statement:
\textit{“The behavioral state is a sufficient statistic for predicting event
probability and determining optimal incentive timing.”}

\section{Recency Dynamics}

Recency captures how long it has been since the user last engaged in the target
behavior. It evolves deterministically:


\[
r_{t+1} =
\begin{cases}
0 & \text{if an event occurs}, \\
r_t + 1 & \text{otherwise}.
\end{cases}
\]



Recency influences event probability through a negative coefficient in the
hazard model, reflecting the empirical observation that the likelihood of
engagement decays with time since last action.

\section{Habit Formation}

Habit strength represents the long-term behavioral momentum of the user. It
increases when events occur and decays during inactivity:


\[
h_{t+1} =
\begin{cases}
\alpha h_t + \gamma & \text{if an event occurs}, \\
\alpha h_t & \text{otherwise},
\end{cases}
\]


where:
\begin{itemize}
    \item $\alpha \in (0,1)$ is the habit retention factor,
    \item $\gamma > 0$ is the reinforcement increment.
\end{itemize}

This formulation matches the habit dynamics described in your uploaded BED
documents:
\textit{“Habit strength evolves over time and influences return probability.”}

\section{Responsiveness}

Responsiveness $\rho_t$ captures how sensitive the user is to incentives. It
represents heterogeneity across individuals and can be:
\begin{itemize}
    \item fixed (segment-based),
    \item estimated (Bayesian updating),
    \item learned (via RL or contextual bandits).
\end{itemize}

In the public-safe BED v3 model, responsiveness is treated as a latent but
slow-moving parameter that modulates incentive effectiveness.

\section{Incentive Decay}

Incentives have diminishing influence over time. BED models this using an
exponential decay function:


\[
\text{EffectiveStrength}_t = I_{t-1} e^{-\lambda \Delta_t},
\]


where:
\begin{itemize}
    \item $I_{t-1}$ is the last incentive magnitude,
    \item $\Delta_t$ is time since that incentive,
    \item $\lambda$ is the decay rate.
\end{itemize}

This aligns with the incentive decay formulation in BED\_MIT.pdf:


\[
\text{EffectiveStrength} = I e^{-\lambda \tau}.
\]



\section{State Update Rules}

The full behavioral state update is:



\[
\begin{aligned}
r_{t+1} &= 
\begin{cases}
0 & \text{if event occurs}, \\
r_t + 1 & \text{otherwise},
\end{cases} \\
\Delta_{t+1} &= 
\begin{cases}
0 & \text{if an incentive is delivered}, \\
\Delta_t + 1 & \text{otherwise},
\end{cases} \\
h_{t+1} &= 
\begin{cases}
\alpha h_t + \gamma & \text{if event occurs}, \\
\alpha h_t & \text{otherwise},
\end{cases} \\
I_t &= a_t,
\end{aligned}
\]


where $a_t$ is the incentive action chosen by the CMDP policy engine.

These update rules match the pseudocode in your uploaded documents (Algorithm 1).

\section{Hazard-Based Event Probability}

BED predicts the probability of an event using a logistic hazard model:


\[
P(\text{event}_t = 1 \mid s_t) =
\sigma\left(
\beta_0
+ \beta_r r_t
+ \beta_h h_t
+ \beta_\rho \rho_t
+ \beta_I I_{t-1} e^{-\lambda \Delta_t}
\right),
\]


where $\sigma$ is the logistic function.

This formulation incorporates:
\begin{itemize}
    \item recency decay,
    \item habit reinforcement,
    \item incentive decay,
    \item responsiveness heterogeneity.
\end{itemize}

It is consistent with the hazard-based response model in BED\_MIT.pdf:


\[
h(\tau \mid I) = h_0(\tau)\exp(\beta I e^{-\lambda \tau}).
\]



\section{Interpretation}

The behavioral model provides a compact, interpretable representation of user
behavior. It enables BED to:
\begin{itemize}
    \item predict event likelihood,
    \item evaluate incentive effectiveness,
    \item update behavioral state in real time,
    \item support CMDP-based policy optimization.
\end{itemize}

\section{Summary}

This chapter formalized the behavioral state model underlying BED v3. The next
chapter presents the CMDP formulation that uses this state representation to
allocate incentives optimally under budget and fairness constraints.

